% =====================================================================
% Chapter 4 — Route A: Multi-View Geometric Fusion
% =====================================================================
\chapter{Route~A: Multi-View Geometric Fusion}
\label{ch:routeA}

This chapter presents the methodology of \routeA, the purely geometric route.
It first justifies \emph{why} Route~A is the appropriate answer in this
project (Section~\ref{sec:why_routeA}), then details the pipeline in its four
stages: pairwise registration (Section~\ref{sec:routea_pairwise}), pose-graph
optimisation (Section~\ref{sec:routea_graph}), fusion (Section~\ref{sec:routea_fusion})
and meshing (Section~\ref{sec:routea_meshing}). Two variants are then described
(Section~\ref{sec:routea_variants}): the view-aware seabed-plane variant and the
sonar-native first-return ray TSDF variant.

\section{Why Route~A?}
\label{sec:why_routeA}

Route~A is the answer to the research question under the following conditions:

\begin{enumerate}
  \item \textbf{Several scans of the same rigid object are available.} The
  \suop{} 3~m acquisition provides 24 curated cases per object, each a
  different placement/orientation of the same physical object. If these scans
  can be aligned into a common frame, each object is observed from many
  effective viewpoints, and the unobserved regions of any single scan are
  covered by the others.
  \item \textbf{The output must be a faithful record of the measurements.}
  Route~A never invents geometry: every point of the output descends from a
  registered sonar return. For safety- and inspection-critical applications
  (e.g.\ estimating whether a propeller is fouled, or whether a net is torn)
  it is preferable to under-cover than to hallucinate, so a measurement-fusion
  route is the defensible choice.
  \item \textbf{No external CAD or learned prior is available or trusted.}
  The \suop{} release contains no CAD models, and the object categories (tyre,
  drum, dummy, net, folding chair) are physically variable, so fitting a rigid
  template would impose assumptions. Route~A requires no prior beyond the
  published object scale used for segmentation.
\end{enumerate}

The \emph{core assumption} of Route~A is that the segmented scans are noisy
partial views of one rigid object, so there exists a set of rigid transforms
that brings them into agreement. The pipeline is therefore a structure-from-scan
problem with unknown poses: registration first, fusion second, meshing third.
Route~A's pipeline is summarised in \figref{fig:routeA_overview}.

\begin{figure}[H]
  \centering
  \includegraphics[width=0.92\textwidth]{fig_routeA_comparison.png}
  \caption{Route~A output for the chair and dummy: fused measured cloud,
  Screened Poisson surface and BPA surface rendered together. The two meshers
  expose the trade-off between global smooth completion (Poisson) and local
  evidence-driven triangulation (BPA).}
  \label{fig:routeA_overview}
\end{figure}

\section{Stage 1: pairwise registration}
\label{sec:routea_pairwise}

\subsection{Preprocessing}
Each segmented cloud is translated so its centroid is at the origin, voxel
downsampled at 15~mm (\texttt{VOXEL\_SIZE}), and given normals estimated with a
hybrid search radius of 45~mm (3~$\times$ voxel). Downsampling serves two
purposes: it reduces the cost of the $\binom{n}{2}$ pairwise evaluations, and
it equalises density between views so that fitness values are comparable.
Normals are required by both the FPFH descriptor and the point-to-plane ICP
refinement.

\subsection{FPFH descriptors}
For each downsampled point, FPFH \citep{rusu2009fast} is computed with a
feature radius of 75~mm (5~$\times$ voxel) and up to 100 neighbours. The
descriptor is a 33-dimensional histogram of relative angle features between
the query point and its neighbours. FPFH is chosen because it is fast, robust
to noise, and adequate for the mostly-smooth objects in \suop{}. Its known
weakness --- confusion between \emph{repeated parts} such as chair legs or
tyre treads --- is exactly why pairwise matches are treated as evidence in a
globally optimised graph rather than accepted as truth.

\subsection{RANSAC global alignment}
For every pair $(i,j)$, an initial alignment is sought with RANSAC feature
matching \citep{fischler1981ransac}: candidate correspondences are sampled from
the FPFH matches, the implied rigid transform is computed, and geometrically
consistent inliers are counted under a correspondence distance of 30~mm
(2~$\times$ voxel). The transform with the most inliers (subject to an
edge-length consistency check and a distance check) is returned. RANSAC is
needed because the \suop{} cases share no common frame and ICP alone cannot
converge from an arbitrary initialisation.

\subsection{Point-to-plane ICP refinement}
The RANSAC transform initialises ICP \citep{besl1992method}, which minimises
the point-to-plane distance between corresponding points,
\begin{equation}
\label{eq:icp}
\min_{T} \sum_{(p,q)\in \mathcal{C}(T)}
\big( (T p - q)\cdot n_q \big)^2,
\end{equation}
where $\mathcal{C}(T)$ is the nearest-neighbour correspondence set under $T$ and
$n_q$ is the normal at the target point. The result is scored by its
\emph{fitness} (fraction of inliers) and \emph{inlier RMSE}. The information
matrix for the edge is computed from the inlier correspondences; it is needed
by the pose-graph optimiser to weight the constraint.

\subsection{Edge acceptance}
Only pairwise registrations whose fitness is at least
\texttt{MIN\_FITNESS\_FOR\_EDGE} $= 0.25$ become pose-graph edges. This
threshold is deliberately conservative: a low-fitness registration typically
aligns only a small (often symmetric) subset of the geometry and would inject
false constraints into the graph. The threshold also encodes the observation
that \suop{} scans at 6 and 10~m ranges simply cannot reach fitness 0.25 (their
fitness is typically $<0.1$ after ICP), which is why only 3~m cases are fused.

\section{Stage 2: pose-graph optimisation}
\label{sec:routea_graph}

\subsection{Maximum-fitness spanning tree}
With $n$ views, the pairwise stage yields up to $\binom{n}{2}$ candidate
edges. A \emph{maximum-fitness spanning tree} is built with Prim's algorithm
\citep{cormen2009introduction} over the accepted edges: starting from a
reference view (index 0), the highest-fitness edge connecting the visited set
to an unvisited view is added until all reachable views are connected. The MST
provides two things: (i)~an \emph{initial pose} for every view (composed along
the tree from the reference), and (ii)~a set of \emph{odometry edges} that are
declared certain in the graph.

\subsection{Loop closures}
Every \emph{other} accepted pairwise edge is a \emph{loop-closure} edge,
declared \emph{uncertain} in the graph. Loop closures are what make multiway
fusion better than a naive pairwise chain: they constrain the accumulated
drift by tying non-neighbouring views together. For 24 views, the MST uses 23
edges and the number of loop closures is typically in the tens to hundreds
(\secref{ch:resultsA}).

\subsection{Global optimisation}
The pose graph is optimised with Open3D's global optimisation
\citep{choi2015robust,zhou2018open3d}: a Levenberg--Marquardt solver minimises
the weighted squared residuals over all edges,
\begin{equation}
\label{eq:posegraph}
\min_{\{T_i\}} \sum_{(i,j)\in E}
\left\lVert \log\!\left( T_{ij}^{-1} T_i^{-1} T_j \right) \right\rVert^2_{\Sigma_{ij}},
\end{equation}
where the norm is the weighted geodesic error on $SE(3)$ and $\Sigma_{ij}$ is
the edge information matrix. An edge-pruning threshold of 0.25 removes
outlier constraints during optimisation. The optimised node poses are the
final per-view object poses in the canonical frame.

\section{Stage 3: fusion}
\label{sec:routea_fusion}

Each original (non-downsampled) segment is transformed by its optimised pose
and accumulated into a single point cloud. The fused cloud is voxel
downsampled at half the pairwise voxel (7.5~mm) to produce the final
\emph{measured} representation of the object. The fused cloud is both an
output in its own right and the input to meshing.

\section{Stage 4: surface reconstruction}
\label{sec:routea_meshing}

Two independent meshers are run on the fused, normal-oriented cloud, exposing
complementary behaviours:

\subsection{Screened Poisson reconstruction}
The cloud's normals are oriented consistently
(\emph{orient\_normals\_consistent\_tangent\_plane}), and Screened Poisson
reconstruction \citep{kazhdan2006poisson,kazhdan2013screened} is solved at
octree depth 8. The output is trimmed by removing the 2\% lowest-density
vertices (a standard guard against extrapolation into unsupported space) and
cleaned: degenerate/duplicate triangles, duplicate vertices and non-manifold
edges are removed, and only the largest connected component is kept.
Poisson's strength is a smooth, closed surface that bridges small sampling
gaps; its limitation is that it \emph{will} close large unsupported gaps too,
inflating sparse regions.

\subsection{Ball-Pivoting Algorithm}
BPA \citep{bernardini1999ball} rolls a ball over the cloud with four radii
($1.5, 2, 3, 4$ $\times$ the mean nearest-neighbour distance), creating a
triangle whenever the ball touches three points without enclosing another.
The same cleaning procedure is applied. BPA's strength is that it preserves
genuine gaps (holes remain where sampling is too sparse); its limitation is
that it cannot bridge gaps at all.

\subsection{Why produce both?}
Producing both meshes is not redundancy: it is the honest way to present the
uncertainty of the fusion. Where Poisson and BPA agree, the surface is well
supported by measurements; where they disagree (Poisson fills, BPA is open),
the surface is extrapolated and should be treated as less certain. This
two-mesh discipline is carried into the results and discussion
(\secref{ch:resultsA}, \secref{ch:discussion}).

\section{Variants}
\label{sec:routea_variants}

Two further Route~A variants address the specific weakness of the generic
pipeline on sparse \suop{} evidence: the generic pipeline registers in full
$SE(3)$, which is fragile when segments contain only a few hundred points.

\subsection{View-aware seabed-plane registration}
The view-aware variant (\texttt{reconstruct\_viewaware.py}) exploits the
physical constraint that the object moves \emph{on the seabed plane}. For each
case it: range-gates the 3~m cloud to the shell band $1.75$--$4.0$~m; removes
the seabed plane (RANSAC, 16~mm threshold); keeps only elevated returns
($3.5$--$48$~cm for the chair, $1.5$--$42$~cm for the tyre) that cannot belong
to the seabed; and clusters at the published object scale (chair diagonal
$0.18$--$0.62$~m; tyre $0.38$--$0.92$~m). Acceptable clusters are registered
with an \emph{SE(2)} transform in the seabed plane --- yaw plus planar
translation --- which also moves the sonar origin into the canonical object
frame, and view-oriented surfels are fused with screened Poisson. The
\emph{acceptance} criterion is explicit: a case is accepted only if its planar
fit has an inlier fraction above threshold and a trimmed RMSE below threshold
(e.g.\ $\le 26$~mm). \tabref{tab:routeA_metrics} in \secref{ch:resultsA} reports
the accepted case counts and the surface-fit metrics.

\subsection{Sonar-native first-return ray TSDF}
\label{sec:routea_tsdf}
The sonar-native variant (\texttt{reconstruct\_suop\_tsdf.py}) deliberately
contains no CAD, primitive or learned prior. It treats every registered return
as the first hit on an acoustic ray anchored at the sonar origin: space
\emph{before} the return is free, and a short band \emph{behind} it is
occupied. These signed-distance observations are fused in a dense TSDF
\citep{curless1996volumetric} (4~mm voxels, 18~mm truncation) and the
zero-crossing is extracted with marching cubes \citep{lorensen1987marching}.
Crucially, the mesh is coloured by \emph{measured-return support} and any
vertex without nearby measured support (beyond a 16~mm support limit) is
removed, so unobserved regions cannot masquerade as accurate geometry. This
variant is the most conservative and the most honest about what a single
acoustic return can prove; its results appear in \secref{sec:resultsA_variants}.

\section{Summary}
\label{sec:routea_summary}

\routeA{} is chosen when several scans of a rigid object exist and the output
must be a faithful fusion of measurements. It estimates poses from geometry
(FPFH/RANSAC/ICP), refines them jointly in a pose graph with loop closures,
fuses the aligned measurements and surfaces the result with two complementary
meshers. Two variants harden the pipeline for sparse sonar evidence: a
seabed-plane-constrained SE(2) registration, and a sonar-native first-return
ray TSDF that reports measurement support explicitly. The next chapter
presents the alternative route, which drops the multi-view requirement
entirely and replaces geometry with a learned prior.
